% GENERATED FILE. Edit canonical lessons, then run npm run generate:books.
% canonical-source-hash: fnv1a64:c3aa5b0a289e63f4
% canonical-lessons: FR-C466-timidite, FR-C466-sociable, FR-C466-reserve, FR-C466-bouder, FR-C466-taquiner

\chapter{Timides et sociables}
\label{ch:fr-timides-et-sociables}
\hlchaptermodality{466}

\begin{chapteropening}
\textbf{By the end of this chapter:} \emph{I can describe how shy or outgoing someone is.}

Complete the last lesson of chapter 466: I can describe how shy or outgoing someone is.
\end{chapteropening}

\section[la timidité]{la timidité — shyness}

\label{lesson:FR-C466-timidite}



\begin{quote}
Before the new one: say the French for a team, then the French for a birthday.
\end{quote}

\subsection*{You'll want to know: la timidité}
\textbf{la timidité} — \textquotedblleft{}shyness\textquotedblright{}.

\begin{quote}
Avec le temps, il a perdu sa timidité.
\end{quote}

\begin{quote}
\emph{With time, he lost his shyness.}
\end{quote}

The adjective \emph{timide} \textquotedblleft{}shy\textquotedblright{} is the everyday word; \emph{la timidité} names the quality.

\begin{cousinweb}
\emph{la timidité} is a learned borrowing from Latin \emph{timiditās} \textquotedblleft{}fearfulness\textquotedblright{} (EtymDB), built on \emph{timidus} \textquotedblleft{}fearful\textquotedblright{}, from \emph{timēre} \textquotedblleft{}to fear\textquotedblright{} (Lewis \& Short).
\end{cousinweb}

\begin{grammarlens}[title={What you've built}]
One more word for telling, explaining and describing at B1.
\end{grammarlens}

\subsection*{Your turn}
\begin{itemize}
\raggedright
  \item \emph{Say it:} \emph{la timidité}
  \item \emph{Say it:} \emph{la timidité}, once more
  \item \emph{Say it:} a sentence of your own with \emph{la timidité}
  \item \emph{Recall:} say the French for a team, then the French for a birthday, then say \emph{la timidité} again
\end{itemize}

\begin{tcolorbox}[breakable,colback=teal!4,colframe=teal!35!black,title={Before you move on}]
What does la timidité mean? (\textquotedblleft{}Shyness\textquotedblright{}.) Say it once more.
\end{tcolorbox}
\section[sociable]{sociable — sociable, outgoing}

\label{lesson:FR-C466-sociable}



\begin{quote}
Before the new one: say the French for a birthday, then the French for shyness.
\end{quote}

\subsection*{You'll want to know: sociable}
\textbf{sociable} — \textquotedblleft{}sociable, outgoing\textquotedblright{}.

\begin{quote}
Ma cousine est très sociable : elle parle à tout le monde.
\end{quote}

\begin{quote}
\emph{My cousin is very sociable: she talks to everyone.}
\end{quote}

The same form for a man and a woman. The opposite is \emph{peu sociable}.

\begin{cousinweb}
\emph{sociable} is a learned borrowing from Latin \emph{sociābilis} \textquotedblleft{}that can be joined together\textquotedblright{}, built on \emph{socius} \textquotedblleft{}a companion, a partner\textquotedblright{} (Lewis \& Short).
\end{cousinweb}

\begin{grammarlens}[title={What you've built}]
One more word for telling, explaining and describing at B1.
\end{grammarlens}

\subsection*{Your turn}
\begin{itemize}
\raggedright
  \item \emph{Say it:} \emph{sociable}
  \item \emph{Say it:} \emph{sociable}, once more
  \item \emph{Say it:} a sentence of your own with \emph{sociable}
  \item \emph{Recall:} say the French for a birthday, then the French for shyness, then say \emph{sociable} again
\end{itemize}

\begin{tcolorbox}[breakable,colback=teal!4,colframe=teal!35!black,title={Before you move on}]
What does sociable mean? (\textquotedblleft{}Sociable, outgoing\textquotedblright{}.) Say it once more.
\end{tcolorbox}
\section[réservé]{réservé — reserved, quiet (of a person); booked}

\label{lesson:FR-C466-reserve}



\begin{quote}
Before the new one: say the French for shyness, then the French for sociable.
\end{quote}

\subsection*{You'll want to know: réservé}
\textbf{réservé} — \textquotedblleft{}reserved, quiet (of a person); booked\textquotedblright{}.

\begin{quote}
Mon frère est un garçon réservé ; il parle peu de lui.
\end{quote}

\begin{quote}
\emph{My brother is a reserved boy; he says little about himself.}
\end{quote}

Of a table or a seat it means booked: \emph{une table réservée}.

\begin{cousinweb}
\emph{réservé} is the past participle of \emph{réserver} \textquotedblleft{}to keep back\textquotedblright{}, which comes from Latin \emph{reservāre} (Lewis \& Short), built on \emph{servāre} \textquotedblleft{}to keep, to watch over\textquotedblright{}.
\end{cousinweb}

\begin{grammarlens}[title={What you've built}]
One more word for telling, explaining and describing at B1.
\end{grammarlens}

\subsection*{Your turn}
\begin{itemize}
\raggedright
  \item \emph{Say it:} \emph{réservé}
  \item \emph{Say it:} \emph{réservé}, once more
  \item \emph{Say it:} a sentence of your own with \emph{réservé}
  \item \emph{Recall:} say the French for shyness, then the French for sociable, then say \emph{réservé} again
\end{itemize}

\begin{tcolorbox}[breakable,colback=teal!4,colframe=teal!35!black,title={Before you move on}]
What does réservé mean? (\textquotedblleft{}Reserved, quiet (of a person); booked\textquotedblright{}.) Say it once more.
\end{tcolorbox}
\section[bouder]{bouder — to sulk}

\label{lesson:FR-C466-bouder}



\begin{quote}
Before the new one: say the French for sociable, then the French for reserved.
\end{quote}

\subsection*{You'll want to know: bouder}
\textbf{bouder} — \textquotedblleft{}to sulk\textquotedblright{}.

\begin{quote}
Le petit garçon boude parce qu'il ne peut pas sortir.
\end{quote}

\begin{quote}
\emph{The little boy is sulking because he cannot go out.}
\end{quote}

\emph{Bouder quelqu'un} means to give someone the cold shoulder.

\begin{cousinweb}
The origin of \emph{bouder} is not settled. It is usually traced to an expressive root \emph{bod-} for something swollen, the picture being pushed-out lips.
\end{cousinweb}

\begin{grammarlens}[title={What you've built}]
One more word for telling, explaining and describing at B1.
\end{grammarlens}

\subsection*{Your turn}
\begin{itemize}
\raggedright
  \item \emph{Say it:} \emph{bouder}
  \item \emph{Say it:} \emph{bouder}, once more
  \item \emph{Say it:} a sentence of your own with \emph{bouder}
  \item \emph{Recall:} say the French for sociable, then the French for reserved, then say \emph{bouder} again
\end{itemize}

\begin{tcolorbox}[breakable,colback=teal!4,colframe=teal!35!black,title={Before you move on}]
What does bouder mean? (\textquotedblleft{}To sulk\textquotedblright{}.) Say it once more.
\end{tcolorbox}
\section[taquiner]{taquiner — to tease (playfully)}

\label{lesson:FR-C466-taquiner}



\begin{quote}
Before the new one: say the French for reserved, then the French for to sulk.
\end{quote}

\subsection*{You'll want to know: taquiner}
\textbf{taquiner} — \textquotedblleft{}to tease (playfully)\textquotedblright{}.

\begin{quote}
Mon oncle aime taquiner les enfants pendant le dîner.
\end{quote}

\begin{quote}
\emph{My uncle likes teasing the children at dinner.}
\end{quote}

Gentle, friendly teasing; stronger mockery is \emph{se moquer de}.

\begin{cousinweb}
\emph{taquiner} is built on the adjective \emph{taquin}, which first meant \textquotedblleft{}stingy\textquotedblright{} and later \textquotedblleft{}fond of teasing\textquotedblright{}. The origin of \emph{taquin} is not settled.
\end{cousinweb}

\begin{grammarlens}[title={What you've built}]
One more word for telling, explaining and describing at B1.
\end{grammarlens}

\subsection*{Your turn}
\begin{itemize}
\raggedright
  \item \emph{Say it:} \emph{taquiner}
  \item \emph{Say it:} \emph{taquiner}, once more
  \item \emph{Say it:} a sentence of your own with \emph{taquiner}
  \item \emph{Recall:} say the French for reserved, then the French for to sulk, then say \emph{taquiner} again
\end{itemize}

\begin{tcolorbox}[breakable,colback=teal!4,colframe=teal!35!black,title={Before you move on}]
What does taquiner mean? (\textquotedblleft{}To tease (playfully)\textquotedblright{}.) Say it once more.
\end{tcolorbox}
